\chapter{揭開希格斯玻色子的奧秘}

\lectureoverview{理解希格斯機制，最好按照各個要素逐步變得必要的順序來看。我們從量子化的運動與場談起，將普通的碗形勢能轉變為墨西哥帽勢能，再利用由此形成的凝聚態，理解費米子和弱玻色子為何能具有質量。接著，我們才將物理希格斯玻色子辨認為凝聚態的徑向振動，並追溯在 2012 揭示它的產生與衰變過程。}

\section{頭條新聞背後的機制}

在 $125\,\mathrm{GeV}$ 附近發現新粒子，確實令人振奮，但振奮還不等於理解。真正值得問的，不是希格斯是否重要，也不是能否在修辭上把它和宇宙起源連結起來。問題其實窄得多：這種現象的具體機制是什麼？標準模型為什麼需要它？

這個名稱也值得稍作限定。彼得·希格斯是指出這項物理機制的數位學者之一，而Fran\c{c}ois恩格勒特與羅伯特·布勞特則各自獨立地提出了核心構想。因此，\emph{Brout--Englert--Higgs 機制}這個名稱更忠實地記錄了歷史，不過下文仍會使用較短的名稱。

這套解釋分成幾個部分。首先需要量子力學的兩個基本事實，以及經典場論的一個概念。接著需要一個最低能量態含有非零場的真空。然後我們可以盤點標準模型中的粒子，看看為什麼該理論的裸形式無法讓相關粒子具有質量。希格斯玻色子要到接近最後才登場。它不是整個機制，而是機制形成後留下的一種特定振動。

\section{兩個量子事實與一個場}

第一個量子事實是量子化。角動量是個清楚的原型，因為它只能取離散的單位，
\begin{equation}
L=n\hbar,
\qquad
n=0,\pm1,\pm2,\ldots .
\end{equation}
半整數角動量也存在，但眼前的論證並不依賴它們。重要的是，連續變化的經典運動會轉變為一系列量子態。

第二個事實是不確定性原理。我們會等到出現一整圈可能真空時，再討論它。屆時，不確定性原理會使我們無法把真空視為一顆靜止在某個精確位置的小小經典球體。

介於這兩個量子事實之間的，是場的經典概念。場會在空間與時間的每一點指定一個量。電場、磁場和重力場都是熟悉的例子。一個區域看起來空無一物，其中仍可能存在場。想像兩片帶電極板遠到看不見，它們之間近乎均勻的電場仍然在當地真實存在。因此，``空無一物的空間''不一定表示``所有場都等於零。''它表示各個場處於可達到的最低能量狀態。

假設場的能量在某個場值處最低，場偏離該值時，能量就會上升。場受到擾動後，會在最低點附近振盪。量子力學將這些振盪量子化，而每一個量子都表現為粒子。整場演講都會用到這個場論對照：
\begin{equation}
\text{particle}
\quad\longleftrightarrow\quad
\text{quantum of a field vibration}.
\end{equation}

\section{場空間中的碗}

對一個實場分量 $\phi$ 而言，最簡單的能量景觀是碗形，
\begin{equation}
V(\phi)=\frac{1}{2}m^{2}\phi^{2}.
\end{equation}
碗底位於 $\phi=0$。碗底附近的小振盪具有確定的頻率，其量子表現為質量為 $m$ 的粒子。

現在將單一場換成兩個分量，$(\phi_1,\phi_2)$。能量只取決於這個內部平面中的半徑，
\begin{equation}
r^{2}=\phi_1^{2}+\phi_2^{2},
\qquad
V=V(r).
\end{equation}
這是\emph{場空間}中的碗，而不是普通空間中的凹陷。旋轉這一對分量 $(\phi_1,\phi_2)$ 不會改變能量。演講中用來說明的帽子，讓這種對稱性變得直觀：繞著帽子的軸旋轉，它看起來依然相同。

\begin{figure}[tbp]
\centering
\includegraphics[width=0.95\textwidth]{lecture_144_figure_02.png}
\caption{普通的對稱碗，以及用來具體呈現場空間旋轉對稱性的實物帽子。}
\lectureframe{00:10:02}
\end{figure}

場一旦開始運動，這個類比就不再只是幾何圖像。將場從中心移開，讓它沿著碗形勢能繞行。這種繞行在場空間中帶有角動量。它不是實驗室中某個物體旋轉的角動量，而是內部角動量。將它量子化後，就得到守恆荷的模型，
\begin{equation}
Q=nq_0,
\qquad
n\in\mathbb{Z}.
\end{equation}
這個方程只是示意。核心的物理概念是，帶電的場激發可以想像成在內部平面中旋轉。正電荷與負電荷對應相反的旋轉方向；沒有旋轉則對應零電荷。

\section{把帽子翻過來}

普通的碗只有一個最低點。希格斯物理從把帽子翻過來開始。此時中心變得不穩定，而一整圈環形谷底具有相同的最低能量。常見的數學形式是
\begin{equation}
V(\phi_1,\phi_2)
=
\lambda
\left(
  \phi_1^{2}+\phi_2^{2}-v^{2}
\right)^{2},
\qquad \lambda>0,
\end{equation}
其極小值滿足
\begin{equation}
\phi_1^{2}+\phi_2^{2}=v^{2}.
\end{equation}
\editorialnote{明確寫出的四次勢能，是將演講中繪出的墨西哥帽幾何圖像標準化為解析形式；它並非轉錄自黑板上的公式。}

\begin{figure}[tbp]
\centering
\includegraphics[width=0.95\textwidth]{lecture_144_figure_03.png}
\caption{倒置帽形勢能的中央峰頂與環形谷底。這張完整的黑板草圖取代了先前只裁出局部、遺漏大部分幾何形狀的圖片。}
\lectureframe{00:13:10}
\end{figure}

真空中的場不再為零，而是必須位於谷底某處，
\begin{equation}
(\phi_1,\phi_2)
=
v(\cos\theta,\sin\theta).
\end{equation}
方程仍保有旋轉對稱性，但任何一個特定的真空都會選定一個角度 $\theta$。這就是自發對稱性破缺：定律具有對稱性，選定的最低能量態卻不具備。

現在有兩種性質迥異的運動。沿徑向向內或向外移動會爬升勢能，因此需要能量。沿谷底繞行只改變 $\theta$；在經典層次上，這種運動不需要額外的勢能。因此，即使帶有大量類似電荷的角向運動，額外勢能也可能非常少。

\begin{figure}[tbp]
\centering
\includegraphics[width=0.95\textwidth]{lecture_144_figure_04.png}
\caption{谷底標記的一點，代表從一整圈等能極小值中選定的一個特定真空。}
\lectureframe{00:16:14}
\end{figure}

\section{不確定性使真空成為凝聚態}

從經典觀點來看，可以把場放在谷底邊緣的一點，令其角速度為零，然後認為問題已經解決。但量子力學不允許如此確定的結論。位置與動量滿足
\begin{equation}
\Delta x\,\Delta p\gtrsim\frac{\hbar}{2}.
\end{equation}
在環形谷底上，對應的說法是：將場的角度局域化，會使與之共軛的內部角動量變得不確定。用演講中的示意性說法表示，
\begin{equation}
\Delta\theta\,\Delta L_{\mathrm{int}}
\gtrsim
\frac{\hbar}{2}.
\end{equation}
\editorialnote{角度算符在技術上有其微妙之處，因此這個關係是演講中不確定性論證的局部半經典表述，並非全域精確的算符恆等式。}

因此，集中在谷底某一點附近的真空，不是具有確定內部電荷的狀態，而是由許多電荷本徵態疊加而成。演講以一個可操作的檢驗來闡明這點：增加或減少一個單位的電荷，幾乎不會改變該狀態。可以用以下示意式理解：
\begin{align}
\ket{\Omega}
&\sim
\sum_n c_n\ket{n},
\\
c_{n+1}
&\simeq
c_n
\quad\text{within the occupied range}.
\end{align}
此時，將 $n$ 增減一，主要只是重新標記原本就已存在的項。這就是稱真空為\emph{凝聚態}的意思：其中含有數目不確定的相關量子，因此多一個或少一個，都不會造成可觀測的宏觀狀態差異。

真實的電弱真空並不是普通電荷的凝聚態。若是如此，空無一物的空間將呈現截然不同的電磁特性。不過，超導體提供了帶電凝聚態的具體例子，因此能作為抽象墨西哥帽圖像與「真空能具有相干量子結構」這項主張之間的物理橋梁。

第一部分到此告一段落。我們看到，最低能量態不必對應零場，而量子不確定性可以使該態對單個量子的增加或移除不敏感。接下來要問的是：哪些粒子需要這樣的背景？

\section{裸標準模型}

基本粒子質量的自然上限尺度是普朗克質量。若將局域能量推高到一定程度，就不能再忽略重力；最終，該物體與其說是普通基本粒子，不如說是黑洞。數值上，
\begin{equation}
m_{\mathrm P}
=
\sqrt{\frac{\hbar c}{G}}
\approx
2.2\times10^{-8}\,\mathrm{kg},
\end{equation}
其質量大致相當於一小粒塵埃。已知基本粒子的質量遠低於此。在演講採用的尺度估算中，它們的質量至多約為普朗克質量的 $10^{-17}$。

這項觀察並不能證明較重的基本粒子不存在。加速器較容易發現能量足以產生、因而較輕的粒子。因此，已知粒子譜在某種程度上受到搜尋設備能力的限制。目前尚無法觸及的範圍內，或許還有許多更重的粒子。

觀測到的粒子分為費米子與玻色子。費米子包括電子、{\proseunicodefont μ}子、{\proseunicodefont τ}子、它們各自的中微子，以及六種夸克味。玻色子包括光子、膠子，以及弱 $W$ 與 $Z$ 玻色子。膠子將夸克束縛成質子和中子，因此在普通原子核的結構中扮演關鍵角色。

\begin{figure}[tbp]
\centering
\includegraphics[width=0.95\textwidth]{lecture_144_figure_05.png}
\caption{演講中第一份標準模型粒子清單，費米子與玻色子兩欄上方畫出了極大的質量範圍。}
\lectureframe{00:24:22}
\end{figure}

這套理論不只是一份清單，也規定了哪些相互作用是允許的。電子可以發射光子，
\begin{equation}
e\longrightarrow e+\gamma.
\end{equation}
夸克帶有電荷與色荷，因此可以發射光子或膠子，
\begin{equation}
q\longrightarrow q+\gamma,
\qquad
q\longrightarrow q+g.
\end{equation}
電子不能發射膠子，而中微子既不發射光子，也不發射膠子。具有適當弱量子數的費米子可以透過 $Z$ 進行相互作用，
\begin{equation}
f\longrightarrow f+Z.
\end{equation}
$W$ 玻色子與帶電流反應也屬於標準模型，但演講暫且略過它們，讓一種中性弱玻色子來承載論證。

\begin{figure}[tbp]
\centering
\includegraphics[width=0.95\textwidth]{lecture_144_figure_06.png}
\caption{光子、膠子和 $Z$ 的放射頂點，旁邊寫著``無質量。''這些過程概括了引入凝聚態之前的裸理論。}
\lectureframe{00:29:20}
\end{figure}

光子沒有質量。在同類型的未破缺規範理論中，膠子與弱規範玻色子也沒有質量。未經修正的電弱理論，同樣不能在保留規範對稱性的同時，直接手動加入費米子質量。理論還缺少一個要素。在引入它之前，我們先從較簡單的情境看看：背景場如何改變粒子的質量，卻不產生摩擦效應。

\section{背景場如何改變質量}

把水分子視為電偶極子：一端略帶正電，另一端略帶負電。在真空中，兩種相反的取向具有相同能量。演講中風趣地將這兩種取向稱為``水''和``威士忌，''只是為了方便把它們當作不同的內部狀態來追蹤。

將分子置於電容器兩極板之間。若其偶極矩為 $\bm d$，電場為 $\bm E$，則交互作用能為
\begin{equation}
U=-\bm d\mathbin{\cdot}\bm E.
\end{equation}
對於兩種分別順向與反向對齊的取向，能量為
\begin{equation}
U_{\parallel}=-dE,
\qquad
U_{\mathrm{anti}}=+dE.
\end{equation}
愛因斯坦關係式 $E=mc^2$ 將能量差轉換為質量差，
\begin{equation}
m_{\mathrm{anti}}-m_{\parallel}
=
\frac{2dE}{c^{2}}.
\end{equation}
分子並沒有受到摩擦。在背景場存在時，只是其中一種內部狀態的能量較高，因此質量也比另一種狀態大。

\begin{figure}[tbp]
\centering
\includegraphics[width=0.95\textwidth]{lecture_144_figure_07.png}
\caption{電容器電場中的兩種相反分子取向。電場造成能量分裂，卻沒有引入拖曳力。}
\lectureframe{00:34:35}
\end{figure}

這個例子排除了兩種容易誤導人的類比。希格斯場不是粒子必須費力穿過的糖漿；有質量的粒子也不是前方堆積了物質、因此減速的推雪車。在均勻的空間中，場不必產生合力。這裡發生的是內部能量的改變，而非摩擦或動量損失。

同一個例子也可以用量子語言重新描述。經典電場是一種相干態；就目前的討論而言，可以將它描述為光子凝聚態。分子可以向這個背景發射光子，也可以從中吸收光子。高度佔據的相干態幾乎不受單一量子的增減影響，因此光子線可以終止於凝聚態。黑板上的叉號表示光子消失於背景之中，並不表示能量不守恆。

\begin{figure}[tbp]
\centering
\includegraphics[width=0.95\textwidth]{lecture_144_figure_08.png}
\caption{將電容器電場改畫為光子凝聚態。偶極子所發射的光子會被背景吸收，而背景幾乎不受額外一個量子影響。}
\lectureframe{00:38:45}
\end{figure}

這就是我們需要的範本。若凝聚態能吸收兩種內部狀態之間的量子數差，就能使這些狀態混合。由此產生的振盪可能表現為質量。

\section{為什麼希格斯並未解釋所有質量}

在套用這個範本之前，得先釐清一個常見誤解：並非所有有質量的物體，其全部質量都來自希格斯場。

將輻射放進完全反射的盒子裡。個別光子沒有質量，但封閉的盒子有靜止系，也儲存著能量。盒子的質量會增加
\begin{equation}
\Delta M=\frac{E_{\mathrm{radiation}}}{c^{2}}.
\end{equation}
因此，即使組成物沒有靜止質量，受限的動能與場能仍能形成質量。

質子是重要的真實例子。它包含處於強相互作用狀態中的夸克、反夸克與膠子。質子的大部分質量來自它們的動能與束縛能，而非單純把少數價夸克由希格斯機制取得的質量加總。在演講強調的估算中，若令輕夸克沒有質量，質子質量只會改變約百分之一。黑洞更鮮明地說明了同一個邏輯：黑洞之所以具有質量，可能是因為形成它的總能量，而不是每一份能量都先由希格斯場轉化為基本粒子的靜止質量。

希格斯機制的作用更為明確：它讓基本費米子與弱規範玻色子得以具有質量，同時維持底層規範理論的一致性。

\section{手徵性：質量是左--右混合}

先看電子。自旋是內禀角動量。對於運動中的粒子，可以比較自旋方向與動量方向。用演講中的幾何說法，右手徵電子與左手徵電子的差異，取決於這兩個方向之間的關係。對無質量粒子而言，螺旋度與手徵性的關係符合這裡所需的簡化說法；對有質量粒子而言，則必須精確區分兩者。

想像自己一再超越一個有質量的電子。在另一個慣性參考系中，電子的動量方向反轉，但自旋方向不變，因此依運動方向定義的表觀手徵性也可能反轉。沒有觀察者能超越嚴格無質量的粒子，因為它永遠以光速運動。這引出演講中的啟發式說法：無質量電子維持在單一手徵態，而有質量電子則持續在左、右手徵態之間混合。

精確的場論說法是，狄拉克質量項會耦合兩個手徵場，
\begin{equation}
\mathcal L_m
=
-m_e\bar e e
=
-m_e
\left(
  \bar e_L e_R
  +
  \bar e_R e_L
\right).
\end{equation}
\editorialnote{狄拉克質量項為演講中的``翻轉率''圖像提供了標準的數學表述。手徵混合的強度與質量成正比；若要把它解讀為實際的時間變化率，則還須考慮粒子的狀態與相對論性時間膨脹。}

若 $m_e=0$，左、右手徵場便各自獨立傳播。若 $m_e\ne0$，方程式就會將兩者耦合起來。就這個精確意義而言，質量衡量左--右手徵態之間的混合強度。

\section{禁止裸質量的守恆定律}

為什麼不一開始就寫入混合項？弱交互作用對 $e_L$ 和 $e_R$ 的作用方式不同。薩斯坎德刻意用一個簡單的名稱來稱呼這項差異：\emph{zilch}。依照演講採用的歸一化方式，
\begin{equation}
z(e_L)=1,
\qquad
z(e_R)=0.
\end{equation}
兩種手徵態具有相同電荷，但弱量子數不同。因此，直接轉變
\begin{equation}
e_L\longrightarrow e_R
\end{equation}
會違反弱交互作用所攜帶的量的守恆。在未破缺的理論中，裸質量項是不允許的。

\editorialnote{``''是教學時採用的簡稱。在標準模型中，左手徵與右手徵費米子屬於不同的 $SU(2)_L\times U(1)_Y$ 表示，因此直接的狄拉克質量項不具規範不變性。演講中的指定值 $1$ 和 $0$ 表達了這項差異，但沒有列出完整的電弱荷表。}

黑板以一連串的過程呈現這項限制。右手徵電子無法轉變為左手徵電子，除非有其他東西帶走兩者之間的 zilch 差異。夸克、{\proseunicodefont μ}子和{\proseunicodefont τ}子也面臨相同限制；補上缺少的場之後，耦合強度則會依粒子種類而異。

\begin{figure}[tbp]
\centering
\includegraphics[width=0.95\textwidth]{lecture_144_figure_09.png}
\caption{黑板上的手徵序列：在綠色叉號處加入一個凝聚態量子，便能讓 zilch 標記在右手徵與左手徵的傳播過程之間改變。}
\lectureframe{00:50:15}
\end{figure}

\section{「``Ziggs''」凝聚態與費米子質量}

引入一個場，讓其量子恰好帶有所缺少的弱量子數。為避免太早與物理希格斯粒子混淆，演講將這個相關量子稱為\emph{Ziggs}。假設 Ziggs 場具有墨西哥帽勢能，且其真空處於凝聚態。這時，左手徵電子可以放出一個 Ziggs 量子而轉變為右手徵電子，
\begin{equation}
e_L\longrightarrow e_R+\text{Ziggs},
\end{equation}
右手徵電子也可以吸收一個 Ziggs 量子而轉變為左手徵電子，
\begin{equation}
e_R+\text{Ziggs}\longrightarrow e_L.
\end{equation}
凝聚態會吸收或提供量子數差異，而自身的宏觀狀態不變。原本在空的未破缺真空中不允許的過程，現在於對稱性破缺的真空中成為可能。

反覆放出與吸收量子，會使 $e_L$ 和 $e_R$ 混合。因此，電子會以有質量粒子的形式傳播。同樣的機制適用於帶電輕子與夸克。各種粒子與凝聚態的耦合強度不同，因此手徵翻轉的強度與質量也會因粒子種類而異。

用現代記法表示，這個規範不變的交互作用稱為湯川耦合。希格斯場取得真空期望值 $v$ 後，會產生
\begin{equation}
\mathcal L_Y
\supset
-\frac{y_f}{\sqrt{2}}
(v+h)\bar f f,
\end{equation}
因此
\begin{equation}
m_f=\frac{y_fv}{\sqrt{2}},
\qquad
g_{hff}=\frac{y_f}{\sqrt{2}}
=
\frac{m_f}{v}.
\end{equation}
\editorialnote{這些公式補上了演講中凝聚態耦合論證所用的標準電弱記法，也明確顯示，同一個耦合如何同時決定費米子質量，以及希格斯衰變至該費米子的過程。}

演講將此描述為手徵對稱性的自發破缺。就這個例子而言，這說明了左--右手徵態之間的關係。更精確地說，希格斯場的真空期望值使電弱規範對稱性破缺至電磁子群，而湯川交互作用隨後成為費米子質量項。

\section{弱玻色子如何獲得質量}

凝聚態也會改變弱規範玻色子的性質。Ziggs 激發可以放出一個 $Z$ 玻色子；反過來，$Z$ 也可以轉變為由凝聚態支持的適當激發，再轉變回來。因此，$Z$ 在一個會與希格斯場角向部分混合的真空中傳播。結果是形成具有質量的向量玻色子，並多出無質量向量玻色子所沒有的縱向極化態。

\begin{figure}[tbp]
\centering
\includegraphics[width=0.95\textwidth]{lecture_144_figure_10.png}
\caption{弱玻色子機制的完整黑板序列：上方仍是手徵鏈，下方則顯示 $Z$ 與凝聚態支持的部分交換量子數。}
\lectureframe{00:53:40}
\end{figure}

這就是以規範玻色子形式呈現的 Brout--Englert--Higgs 現象。在完整的標準模型中，$W$玻色子與$Z$玻色子都會獲得質量，而原先的中性規範場中，有一個組合仍成為無質量的光子。假如真空破缺的是一般的電磁規範對稱性，光子就會變得有質量，正如光子在超導體內部實際上呈現的情形。如此一來，真空中的長程電磁行為就會截然不同。

在物理希格斯粒子被單獨辨認出來之前，弱玻色子早已為這項機制提供了有力證據。它們的質量和交互作用顯示，電弱對稱性確實以所需的方式破缺。因此，演講可以用較寬鬆的操作性說法指出，``Ziggs''部分早已為人所知；實驗上尚待找到的是獨立的徑向振動。

\editorialnote{演講中即興提及約在 1980 年於 SLAC 發現的事件，實際上概括了好幾件事。中性弱流於 1973 年被觀測到；SLAC 在 1970 年代後期報告了宇稱不守恆的電子散射；而 $W$ 與 $Z$ 玻色子則於 1983 年在 CERN 首度直接發現。這裡的物理重點，是電弱對稱性破缺早已獲得實驗確認。}

\section{希格斯玻色子的本質}

回到墨西哥帽。沿谷底運動會改變角度，但不改變半徑。這個角向部分提供了被有質量弱玻色子吸收的模式。物理希格斯玻色子對應另一種運動：沿著勢能的斜面向內、向外振盪。

將場示意性地寫成
\begin{equation}
\Phi(x)
=
\frac{v+h(x)}{\sqrt{2}}
\exp\!\left(
  \frac{i\theta(x)}{v}
\right).
\end{equation}
變數 $\theta$ 描述沿谷底的運動；變數 $h$ 則改變半徑。在規範理論中，適當的角向自由度會成為有質量弱玻色子的縱向分量，而 $h$ 則保留下來，成為可以傳播的純量粒子：希格斯玻色子。

\begin{figure}[tbp]
\centering
\includegraphics[width=0.95\textwidth]{lecture_144_figure_11.png}
\caption{黑板上再次畫出兩種勢能，並在墨西哥帽的谷底標示徑向位移。這種向內、向外的振動就是物理希格斯模式。}
\lectureframe{00:56:30}
\end{figure}

也可以用多粒子圖像來理解。真空處於凝聚態，徑向振盪會改變凝聚態的密度。因此，希格斯玻色子就像這個量子介質中的壓縮波。只要不因此認為波必然穩定，這個類比就很有用。普通聲波會衰減，希格斯激發也可能衰變。

這項區分說明了為什麼發現希格斯玻色子是補上最後一塊拼圖，而不是首次證實這項機制。破缺的真空與有質量的弱玻色子早已間接顯現。實驗上還必須產生徑向密度波，並透過其衰變產物辨認它。

\section{衰變、產生與重費米子的優勢}

希格斯與某種費米子的耦合強度正比於該費米子的質量。希格斯可以衰變成一對費米子與反費米子，
\begin{equation}
H\longrightarrow f+\bar f,
\end{equation}
同一張費曼圖反向解讀時，則表示粒子生成過程，
\begin{equation}
f+\bar f\longrightarrow H.
\end{equation}
因此，電子與正電子湮滅而產生希格斯，在原理上是可行的。然而，早期對撞機的能量不足以產生希格斯；後來的機器則面臨更根本的問題：電子的湯川耦合極小。輕夸克的湮滅過程也同樣微弱。

\begin{figure}[tbp]
\centering
\includegraphics[width=0.95\textwidth]{lecture_144_figure_12.png}
\caption{黑板上畫出的希格斯衰變頂點。反向解讀箭頭，便是費米子與反費米子融合生成希格斯。}
\lectureframe{00:59:40}
\end{figure}

費米子越重，湯川耦合就越大。這使頂夸克格外突出：它的質量約為$173\,\mathrm{GeV}$，相當於約$180$個質子的質量。講座為了方便比較，採用約$170$個質子質量這個整數近似值。然而，質量接近$125\,\mathrm{GeV}$的希格斯太輕，無法衰變成一對在殼的$t\bar t$，
\begin{equation}
m_H<2m_t.
\end{equation}
也無法製成實用的頂夸克束流：頂夸克還來不及強子化成能長時間存在的束流粒子，就已經衰變。

量子力學提供了克服這兩個障礙的方法。兩個膠子可以經由虛頂夸克圈發生漲落，最後產生一個希格斯，
\begin{equation}
g+g\longrightarrow H.
\end{equation}
虛頂夸克線不必滿足真實頂夸克對產生所受的在殼閾值限制。質子中充滿膠子，因此大型強子對撞機的高能質子—質子碰撞能大量提供所需的膠子對。

\begin{figure}[tbp]
\centering
\includegraphics[width=0.95\textwidth]{lecture_144_figure_13.png}
\caption{透過虛頂夸克圈的膠子融合。取自碰撞質子的兩個膠子在圈中相遇，並產生希格斯。}
\lectureframe{01:06:35}
\end{figure}

因此，儘管質子碎片相當複雜，質子對撞機仍成為關鍵的希格斯工廠。質子本身並非以單一基本粒子的身分與希格斯耦合；而是質子中的膠子組分透過一個由頂夸克的強耦合主導的圈，產生希格斯。

\section{2012粒子與雙光子檢驗}

2012年報告的粒子，其質量接近
\begin{equation}
m_H\simeq125\,\mathrm{GeV}.
\end{equation}
依照講座中便於比較的數值，這相當於約$127$個質子的質量。它的發現補齊了最小標準模型所要求的粒子清單。接下來的問題是，它的定量性質是否完全符合標準模型的預測。

頂夸克圈的道理也可以反過來看。希格斯與圈中的虛帶電粒子耦合，而這個圈可以放出兩個光子，
\begin{equation}
H\longrightarrow\gamma+\gamma.
\end{equation}
這種衰變很罕見，但在強子對撞機中，光子是特別容易辨認的粒子。因此，即使較常見的衰變在實驗上更難分析，仍可從雙光子不變質量分布中的窄峰辨認出希格斯。

\begin{figure}[tbp]
\centering
\includegraphics[width=0.95\textwidth]{lecture_144_figure_14.png}
\caption{黑板右側的雙光子三角圖，對應左側先前畫出的膠子融合圈。這兩個過程從相反方向檢驗同一套圈圖層級的物理原理。}
\lectureframe{01:09:00}
\end{figure}

講座當時，報告中的雙光子衰變率似乎約為最簡單標準模型預測值的一點五倍，但顯著性只有約兩個標準差。這樣的漲落本身並不足以構成證據。然而，若它持續存在，就可能表示圈中還有另一種帶電粒子，或許與超對稱或其他標準模型延伸理論有關。

這段話反映的是2012當時的歷史情境。它記錄了當時實驗物理學家被告知接下來該關注什麼，並非在宣稱目前仍有異常。希格斯的發現本身大致上確認了一個早已獲弱玻色子資料支持的理論框架。更具判別力的科學檢驗在於各種衰變率：這個新標量粒子與各粒子的耦合，是否呈現希格斯機制所預測的質量依賴模式？

\section{聽眾提問}

\begin{classroomqa}
\audiencequestion{是什麼決定某種費米子在左、右手態之間振盪的速度比另一種快，因而具有不同的質量？\lecturetimestamp{01:12:00}}
\lecturerresponse{每一種費米子都以自己的耦合強度與希格斯凝聚體作用。耦合越強，左--右手態的混合就越強，質量也越大。同一個耦合也決定了物理希格斯衰變成該種費米子的強度，這就是為什麼較重的費米子通常耦合較強。標準模型以湯川耦合參數描述這些耦合，卻沒有解釋它們為何呈現如此顯著的階層。說電子的耦合較弱、頂夸克的耦合較強，描述是正確的，但還稱不上更深入的推導。}
\end{classroomqa}

\begin{classroomqa}
\audiencequestion{我們知道真空中希格斯場的大小嗎？\lecturetimestamp{01:13:36}}
\lecturerresponse{知道。依照這裡採用的慣例，真空期望值約為$240\,\mathrm{GeV}$；較精確的標準值約為$246\,\mathrm{GeV}$。可以把它想成場從原點沿徑向的位移，也可以把它想成凝聚體的密度。物理希格斯量子就是這個值的小幅徑向振盪；等價地說，它是凝聚體局部密度的漲落。}
\end{classroomqa}

\begin{classroomqa}
\audiencequestion{希格斯場非零的真空值，是否表示空間真空具有能量密度？\lecturetimestamp{01:14:12}}
\lecturerresponse{是。希格斯勢會對真空能有所貢獻，量子場論中的許多其他部分也一樣。觀測到的淨真空能與這些貢獻的自然尺度相比，微小得驚人，因此各項貢獻必須幾乎彼此抵消。目前沒有公認的理論能解釋這種抵消。這就是宇宙學常數問題，也是基礎物理學中最深刻的未解問題之一。}
\end{classroomqa}

\section{這套建構的成果}

現在可以從頭到尾理解這套推理，而不必借助黏滯以太。量子化的內部旋轉提供了理解電荷的模型。把碗形勢能改成墨西哥帽形，場便具有非零的真空值。不確定性原理使原本精確選定的經典真空轉變成凝聚體，不再受個別量子的影響。接著，凝聚體可以承載左、右手費米子之間的弱同位旋量子數差，讓兩者得以混合，因而獲得質量。凝聚體也提供縱向模式，使$W$和$Z$能以有質量向量玻色子的形式傳播，同時讓光子保持無質量。

重新安排之後留下的徑向振盪，就是物理希格斯玻色子。它不是``質量本身，''也不負責質子的大部分質量。它是對稱破缺場中可觀測到的密度波。它的耦合反映了該場所產生的粒子質量，這就是為什麼重粒子在它的產生與衰變過程中占主導地位。因此，膠子融合圈、雙光子衰變道，以及聽眾最後提出的真空能問題，並非互不相關的趣聞。它們分別是同一個非平凡真空的後果、檢驗，以及尚待釐清的含意。
